\begin{textsample}{POS Dim 3 – sp – Score 18.00 – sp/sp\_inf\_funcao\_social.txt}  \label{ex:f3_pos_203}
Função social da \textbf{instituição} de Educação \textbf{Infantil} A \textbf{instituição} de Educação \textbf{Infantil}, responsável pela primeira etapa de Educação Básica, visa a atender as especificidades da \textbf{criança} \textbf{pequena} sem, contudo, ser preparação para o Ensino Fundamental.
Assim, contrapondo -se à ideia de preparatória, essa etapa exige priorizar as interações e as \textbf{brincadeiras} como eixos estruturantes para a organização de tempos e espaços, de modo a garantir experiências ricas para a aprendizagem, o que não combina com a proposição de atividades estanques, fragmentadas.
Uma \textbf{instituição} de Educação \textbf{Infantil} que prioriza as interações e a \textbf{brincadeira} tem a prática de \textbf{ouvir} as \textbf{crianças}, por exemplo, sobre como podem ser dispostos os \textbf{brinquedos} no \textbf{parque}, como deve ser organizada a \textbf{biblioteca}, os espaços, a adequação e disposição das mobílias.
Assim, abre espaços e possibilidades para que as \textbf{crianças} participem nas diversas decisões, inclusive no planejamento da gestão da escola e das atividades propostas pelo educador ( BRASIL, 2017 ).
É importante destacar que a atenção ao que a \textbf{criança} fala não se encerra na linguagem verbal, mas esta deve considerar as sutilezas das formas de comunicação dos \textbf{bebês} e das \textbf{crianças}, como afirma Loris Malaguzzi, revelado no \textbf{livro} As cem linguagens da \textbf{criança} : “ [... ] A \textbf{criança} tem cem \textbf{mãos}, cem pensamentos, cem modos de pensar, de jogar e de falar [... ] ” ( EDWARDS, et al, 1999, p.5 ).
Deste modo, cabe ao professor \textbf{ouvir} não apenas com ouvidos, mas com olhar responsivo, observando as expressões de cada \textbf{criança}, \textbf{acolhendo} e inferindo as necessidades e interesses dela a partir do que observa.
As \textbf{crianças} precisam ser consideradas também quanto à disposição e às quantidades de \textbf{mobiliário} da sala, levando em conta suas especificidades e a necessidade de movimentar -se, explorar diferentes espaços, criar cenários, \textbf{brincar} junto com outras \textbf{crianças}.
Em vista disso, a BNCC, como política pública, elege como núcleo da nova Educação \textbf{Infantil} as \textbf{crianças} e suas experiências, assegurando -lhes o direito de aprender e se desenvolver.
O diálogo da Educação \textbf{Infantil} com outros setores Pensar o desenvolvimento integral da \textbf{criança} requer considerá -la nos diferentes contextos sociais.
A indissociabilidade do \textbf{cuidar} e do educar demanda diversas ações das \textbf{instituições} públicas, de maneira especial, dos equipamentos públicos da comunidade onde a escola está inserida ; e prevê uma articulação orquestrada, na qual diferentes agentes tecem, por meio das suas atuações, uma rede de proteção à \textbf{infância}.
É desejável que a ação intersetorial esteja explicitada no projeto político pedagógico da escola considerando o contexto local, uma vez que, conforme afirmado nos Parâmetros Nacionais de Qualidade da Educação \textbf{Infantil} ( PNQEI, 2006 ), “ a proteção integral das \textbf{crianças} extrapola as funções educativas e de \textbf{cuidado} e deve ser articulada por meio de ações que integrem as políticas públicas intersetoriais ”.

% matched lemmas: acolher, bebê, biblioteca, brincadeira, brincar, brinquedo, criança, cuidado, cuidar, infantil, infância, instituição, livro, mobiliário, mão, ouvir, parque, pequeno
\end{textsample}
